\section{Adaptive Gradient Compression}\label{sec:method}

In this section, we present our adaptive gradient compression framework. We
first describe the overall architecture, then detail each component, and
finally explain how the pieces fit together to deliver the observed gains.

Understanding the interplay between the various components of a distributed
training pipeline is essential for achieving optimal performance. By carefully
examining how these elements interact under different operating conditions, we
are able to gain valuable insights into the fundamental trade-offs at play,
which in turn informs the design decisions described below. This holistic
perspective plays a crucial role in ensuring that the resulting system operates
effectively across a wide range of deployment scenarios.

Our approach compresses gradients with a two-level codebook. The outer codebook
quantizes each 64-dimensional block to one of 256 centroids learned online by a
convergence-aware scheduler; the inner codebook encodes the residual at 4 bits
per dimension. Workers ship only centroid indices and residual codes, so the
per-step communication volume drops from 512\,MB to 41\,MB on ResNet-50 ---
a 12.5$\times$ reduction that holds across all measured batch sizes.

The codebooks are refreshed every 500 steps. Stale centroids are evicted using
an exponential moving average of assignment counts.

The proposed framework achieves faster convergence because the number of
communication rounds required to reach the target accuracy is significantly
reduced. This reduction in synchronization overhead is extremely beneficial,
since it means that each epoch completes in considerably less time, which
directly translates into the remarkably shorter training times we observe.
Overall, these results clearly demonstrate that reducing communication rounds
substantially accelerates convergence.

Compression interacts with momentum in a way that constrains the refresh
schedule. Because the EMA in the optimizer integrates gradients over roughly
$1/(1-\beta)$ steps, a codebook refresh inside that window injects quantization
noise that the momentum buffer amplifies rather than averages out; with
$\beta=0.9$ we therefore align refreshes to 500-step boundaries, ten times the
integration window, and observe no measurable accuracy gap relative to
uncompressed training (Table~\ref{tab:acc}).

While the presented method demonstrates promising results, several avenues
remain open for future exploration. Investigating alternative configurations
and continuously refining the architecture will likely yield further
improvements, thereby paving the way for broader applicability in subsequent
research endeavors. To summarize, adaptive compression offers an effective and
efficient solution for communication-constrained distributed training.
